\chapter{The Visual Front End}
\label{ch:frontend}

\begin{objectives}
\begin{itemize}
  \item Understand the structure tensor and why it defines both good corners and solvable tracking.
  \item Derive Lucas--Kanade optical flow from brightness constancy, including the aperture problem and pyramids.
  \item Build a robust tracking pipeline: bucketing, forward-backward checks, IMU-aided prediction and RANSAC.
  \item Manage feature tracks through their life cycle with bounded memory.
\end{itemize}
\end{objectives}

The front end turns a stream of images into \emph{feature tracks}: for each tracked scene point, the list of its
undistorted normalized coordinates in the images where it was seen, with the image timestamps. Everything the MSCKF
knows about the scene comes from these tracks, so their quality bounds the quality of the estimator.

\begin{figure}[h]
\centering
\begin{tikzpicture}[node distance=4mm and 5mm, every node/.style={block, font=\footnotesize, minimum width=19mm}]
  \node (img) {image $k$};
  \node (pyr) [right=of img] {build\\pyramid};
  \node (lk) [right=of pyr] {track with\\pyramidal LK};
  \node (fb) [right=of lk] {forward-\\backward check};
  \node (und) [below=of fb] {undistort\\points};
  \node (ran) [left=of und] {RANSAC\\(E or 2-pt)};
  \node (det) [left=of ran] {detect new\\corners (grid)};
  \node (out) [left=of det, fill=white] {tracks to\\estimator};
  \draw[->] (img) -- (pyr); \draw[->] (pyr) -- (lk); \draw[->] (lk) -- (fb);
  \draw[->] (fb) -- (und); \draw[->] (und) -- (ran); \draw[->] (ran) -- (det); \draw[->] (det) -- (out);
\end{tikzpicture}
\caption{A KLT-based front end, as used by MSCKF-style estimators.}
\label{fig:frontend-pipeline}
\end{figure}

% ------------------------------------------------------------------------------
\section{What makes a good feature}

Consider shifting a small window $W$ of the image $I$ by a displacement $d$. The sum of squared differences is,
to second order,
\begin{equation}
  E(d) = \sum_{x\in W}\big(I(x + d) - I(x)\big)^2 \approx \sum_{x\in W}\big(\nabla I(x)\T d\big)^2 = d\T M\,d,
  \qquad
  M = \sum_{x\in W}\begin{bmatrix} I_x^2 & I_xI_y \\ I_xI_y & I_y^2 \end{bmatrix},
  \label{eq:structure-tensor}
\end{equation}
where $I_x$, $I_y$ are the image gradients. $M$ is the \emph{structure tensor} (second-moment matrix) of the window. Its
eigenvalues $\lambda_1 \ge \lambda_2 \ge 0$ describe how distinctive the window is.
\begin{itemize}
  \item Both small: the window is \emph{flat}; it looks the same after any small shift.
  \item One large, one small: an \emph{edge}; shifting along the edge changes nothing.
  \item Both large: a \emph{corner}; every shift changes the window, so its position is well defined in 2D.
\end{itemize}
The Harris detector~\cite{harris1988corner} scores $\det M - k(\tr M)^2$ to avoid computing eigenvalues; the Shi--Tomasi
detector~\cite{shi1994good} scores $\lambda_2 = \min(\lambda_1, \lambda_2)$ directly (\cref{fig:structure-tensor}). The two behave similarly; Shi--Tomasi has a
deeper justification, which the next section reveals.

\begin{figure}[t]
  \centering
  \includegraphics[width=\textwidth]{fig_structure_tensor.pdf}
  \caption{Corner scores over the plane of structure-tensor eigenvalues, with a threshold contour (orange). Both
  scores accept only windows with two large eigenvalues. Shi--Tomasi's $\min(\lambda_1, \lambda_2)$ is exactly the quantity that
  makes the Lucas--Kanade system well conditioned.}
  \label{fig:structure-tensor}
\end{figure}

% ------------------------------------------------------------------------------
\section{Lucas--Kanade optical flow}

\subsection{Brightness constancy and the aperture problem}

Assume that a scene point keeps its brightness between consecutive frames:
$I_{k+1}(x + d) = I_k(x)$. Linearizing in $d$,
\begin{equation}
  I_{k+1}(x) + \nabla I\T d \approx I_k(x)
  \quad\Longleftrightarrow\quad
  I_x\,d_x + I_y\,d_y + I_t = 0, \qquad I_t = I_{k+1}(x) - I_k(x).
  \label{eq:ofc}
\end{equation}
This optical-flow constraint is one equation in two unknowns. It determines only the component of $d$ along the
gradient: motion along an edge is invisible. This is the \emph{aperture problem}, the same phenomenon seen through a
small hole when a long line slides past.

\subsection{The Lucas--Kanade solution}

Lucas and Kanade~\cite{lucas1981iterative} assume the displacement is constant over a small window and solve
\eqref{eq:ofc} for all its pixels in the least-squares sense:
\begin{equation}
  \min_d \sum_{x\in W}\big(\nabla I(x)\T d + I_t(x)\big)^2
  \quad\Longrightarrow\quad
  M\,d = -\sum_{x\in W} I_t(x)\nabla I(x),
  \label{eq:lk}
\end{equation}
with the structure tensor $M$ of \eqref{eq:structure-tensor} as the normal matrix. The system is solvable and well
conditioned exactly when both eigenvalues of $M$ are large compared with the image noise. Shi and Tomasi's insight
was to select features by the very criterion that makes them trackable: a ``good feature to track'' is a window
with a large $\lambda_{\min}(M)$.

Because \eqref{eq:lk} rests on a linearization, it is iterated. Warp the second image by the current estimate, recompute
$I_t$, solve for an increment, and repeat until the increment is a small fraction of a pixel. This is Gauss--Newton on
the sum of squared differences~\cite{baker2004lk}. Sub-pixel accuracy comes from bilinear interpolation of the
warped image.

\subsection{Pyramids}

The linearization is valid only for displacements smaller than the scale over which the image gradient is roughly
constant, about the window radius. Fast motion violates this. The remedy is a Gaussian \emph{pyramid}: the image is
repeatedly blurred and down-sampled by two (\cref{fig:pyramid}). Tracking starts at the coarsest level, where a
40-pixel motion becomes a 5-pixel motion, and the estimate is doubled and refined at each finer
level~\cite{bouguet2001pyramidal}. With $L$ levels above the base, the trackable displacement grows by about $2^L$.

\begin{figure}[t]
\centering
\begin{tikzpicture}[scale=0.9]
  \foreach \l/\s/\y in {0/3.2/0, 1/1.6/2.0, 2/0.8/3.2, 3/0.4/3.95} {
    \draw[fill=primarylight, draw=primary] (-\s, \y) -- (\s, \y) -- ({\s*0.75}, {\y + 0.35}) -- ({-\s*0.75}, {\y + 0.35}) -- cycle;
    \node[font=\footnotesize, right] at ({\s + 0.2}, {\y + 0.15}) {level \l: $\frac{1}{2^{\l}}$ resolution};
  }
  \draw[->, plotorange, thick] (-0.3,4.2) -- (-0.15,3.55) -- (0.2,2.4) -- (0.9,0.25);
  \node[plotorange, font=\footnotesize, align=right, left] at (-0.5,2.6) {coarse estimate,\\refined at each level};
\end{tikzpicture}
\caption{An image pyramid. A displacement that is too large to linearize at full resolution is small at the top
level; each level doubles the estimate and refines it.}
\label{fig:pyramid}
\end{figure}

% ------------------------------------------------------------------------------
\section{Making tracking robust}

Optical flow optimizes photometric similarity, not physical correctness. It will track the shadow of a moving object,
a specular highlight, or a T-junction formed by two objects at different depths. Several layers of checks remove such
tracks.
\begin{enumerate}
  \item \textbf{Forward-backward consistency.} Track each point from frame $k$ to $k+1$, then back to $k$. A correct track
        returns to its starting point; reject it if the round-trip error exceeds a fraction of a pixel.
  \item \textbf{IMU-aided prediction.} The gyroscope predicts the rotation between frames. Predicting each feature's new
        position as $\pi(R_{c_{k+1}c_k}\bar x_k)$ (exact for distant points) gives LK a much better starting point and
        allows smaller windows and fewer pyramid levels.
  \item \textbf{Geometric verification.} Reject tracks that violate the epipolar constraint, with essential-matrix or
        two-point RANSAC, or the rotation-only constraint under low parallax (\cref{ch:multiview}).
  \item \textbf{Statistical gating in the filter.} Tracks that pass all of the above are still gated by the $\chi^2$ test
        of the MSCKF update (\cref{ch:msckfupdate}).
\end{enumerate}

\subsection{Feature distribution}

Asking a detector for the 200 strongest corners typically returns clusters on the most textured object in view.
Clustered features constrain the motion poorly: features in a small patch of the image make rotation and
translation hard to separate. The standard remedy is \emph{grid bucketing}: divide the image into cells (for example
$8\times6$) and cap the number of features per cell, plus a minimum distance between features. New features are
detected only in cells that have fallen below their quota, using a mask around existing tracks. A uniform distribution
of bearings makes the stacked measurement Jacobian better conditioned, which is exactly what the estimator needs.

% ------------------------------------------------------------------------------
\section{The feature life cycle}

\begin{figure}[h]
\centering
\begin{tikzpicture}[node distance=13mm and 34mm, every node/.style={font=\small}]
  \node[block] (new) {NEW};
  \node[block, right=of new] (trk) {TRACKED};
  \node[block, right=of trk] (lost) {LOST};
  \node[block, below=of trk, fill=white] (used) {consumed by update\\or discarded};
  \draw[->, thick] (new) -- node[above, font=\footnotesize]{tracked} (trk);
  \draw[->, thick] (trk) edge[loop above] node[above, font=\footnotesize]{tracked again} (trk);
  \draw[->, thick] (trk) -- node[above, font=\footnotesize]{track fails} (lost);
  \draw[->, thick] (lost) |- (used);
  \draw[->, thick] (trk) -- node[right, font=\footnotesize, align=left]{oldest observation's\\clone leaving the window} (used);
  \draw[->, thick] (new) |- node[left, font=\footnotesize, pos=0.25]{rejected} (used);
\end{tikzpicture}
\caption{Feature life cycle. A track is handed to the MSCKF update when it is lost, or when its oldest observation is
about to leave the sliding window (AGENTS.md F5). Its observations are used once and then discarded.}
\label{fig:lifecycle}
\end{figure}

Each feature has a persistent ID and a history of observations: (clone index or image timestamp, normalized
coordinates). Since the estimator can only use observations made at times that still have a clone in the window,
the history never needs to exceed \code{kMaxClones} entries. That makes it possible to store all tracks in a
fixed-capacity pool with fixed-capacity histories, so the steady state of the front end performs no allocation for
its bookkeeping (Milestone 9).

\begin{pitfall}[using a measurement twice]
Once a track's observations have been used in an MSCKF update, they must never be used again, not even partially.
Re-using them double counts the same information and makes the filter overconfident. If a track continues after being
used (because its oldest clone was marginalized), its subsequent observations start a fresh track.
\end{pitfall}

% ------------------------------------------------------------------------------
\section{Engineering the front end}

The front end is the one component allowed to allocate (AGENTS.md E4), because OpenCV allocates internally, but it
must reuse its buffers. In practice: keep the previous pyramid and swap it with the new one instead of rebuilding
storage; preallocate point vectors with \code{reserve}; keep the \code{cv::Mat} objects as members so that their memory is
reused when sizes do not change; and \emph{measure} the remaining allocations with the allocation guard rather than
claiming they are zero. Feature tracking is typically the most expensive part of the pipeline in time, too, so bucketing
parameters, window size and pyramid depth are the main performance knobs.

Descriptor-based front ends (ORB, BRIEF) match features across frames by appearance instead of tracking them. They
handle large viewpoint changes and re-detection after occlusion, which matters for loop closure, but for frame-to-frame
VIO at 20~Hz KLT tracking is cheaper and more accurate. Direct methods skip features altogether and use photometric
error in the estimator itself.

\begin{projectlink}[Milestones 8 and 9]
M8 builds detection, bucketing, pyramidal LK with forward-backward checks, persistent IDs and point undistortion, and
tests them on synthetic images with known sub-pixel motion. M9 adds the life-cycle manager with a fixed-capacity pool
and essential-matrix RANSAC in normalized coordinates, including the documented pure-rotation behavior.
\end{projectlink}

\begin{takeaways}
\begin{itemize}
  \item The structure tensor $M$ measures how well a window's position is defined; corners have two large eigenvalues.
  \item Lucas--Kanade solves $Md = -\sum I_t\nabla I$: good corners are exactly the trackable windows.
  \item The aperture problem: a single pixel constrains only the flow along the gradient.
  \item Pyramids extend the trackable displacement; gyro prediction makes tracking cheaper and more robust.
  \item Bucketing spreads features for good geometry; forward-backward checks and RANSAC remove bad tracks.
  \item Tracks are consumed exactly once and their histories are bounded by the window size.
\end{itemize}
\end{takeaways}

\begin{exercises}
  \item Compute the structure tensor of an ideal step edge and of an ideal corner, and verify the eigenvalue picture.
  \item Implement single-level LK for a synthetic textured image translated by $(0.3, -0.7)$ pixels. How many iterations
        until convergence to $0.01$ px? What happens with a translation of 8 pixels, and how many pyramid levels fix it?
  \item Show that the forward-backward error of a correct track is due only to noise and interpolation, and estimate a
        sensible threshold.
  \item Implement grid bucketing and plot the feature distribution with and without it on a EuRoC image.
  \item Design the data structure for the feature pool with fixed capacity and per-feature histories bounded by
        \code{kMaxClones}. What happens when the pool is full?
\end{exercises}
